\documentclass[12pt]{article}

\usepackage{amssymb,amsmath,amsthm}

% gaya teorema (saya ingin semuanya menggunakan pencacah yang sama)
\theoremstyle{plain}
\newtheorem{theorem}{Teorema}[section]
\newtheorem{lemma}[theorem]{Lema}
\newtheorem{proposition}[theorem]{Proposisi}
\newtheorem{corollary}[theorem]{Korolari}
\newtheorem{conjecture}[theorem]{Konjektur}
\newtheorem{exercise}[theorem]{Latihan}
\theoremstyle{definition}
\newtheorem{definition}[theorem]{Definisi}
\newtheorem{example}[theorem]{Contoh}
\theoremstyle{remark}
\newtheorem{remark}[theorem]{Catatan}

% beberapa pengaturan penomoran
\numberwithin{equation}{section}

\usepackage{fancyhdr}
\usepackage{lastpage}
\setlength{\headheight}{15.2pt}
\renewcommand{\footrulewidth}{0.4pt}% nilai bawaan adalah 0pt
\setlength{\footskip}{30pt}
\pagestyle{fancy}

\makeatletter
\let\ps@plain\ps@fancy 
\makeatother

\lhead{MATH 742}
\chead{Teori Karakter}
\rhead{Musim Semi 2023}
\lfoot{Terakhir Direvisi: \today}
\cfoot{}
\rfoot{\thepage\ dari \pageref{LastPage} }

\begin{document}

%%%%%%%%%%%%%%%%%%
%%%%%%%%%%%%%%%%%%
\title{Teori Karakter}
\author{Alexander Duncan}

\maketitle

{\bf
Di seluruh dokumen ini, kita menggunakan asumsi tetap bahwa $G$
adalah grup hingga dan $k$ adalah medan dengan karakteristik yang relatif prima terhadap
orde $G$.}
Dengan demikian, kita dapat menggunakan teorema Maschke:
setiap representasi berdimensi hingga tereduksi sepenuhnya.
Jadi, setiap representasi tak terurai juga tak tereduksi.

Teori karakter merupakan inti teori representasi grup hingga.
Sebagian besar materi di bawah ini dapat ditemukan, misalnya, dalam
\cite[\S{14--15}]{AlperinBell},
\cite[\S{18.3}]{DF},
\cite[\S{4}]{Etingof},
\cite[\S{2}]{FultonHarris},
\cite[\S{XVIII.1--5}]{Lang}, atau
\cite[\S{2}]{Serre}.

\section{Gelanggang Representasi}

Misalkan $R^+_k(G)$ adalah himpunan kelas isomorfisme
representasi berdimensi hingga dari $G$ atas $k$.
Untuk representasi $V$ dan $W$, misalkan $[V]$ dan $[W]$ menyatakan
bayangannya dalam $R^+_k(G)$.
Kita mendefinisikan penjumlahan pada $R^+_k(G)$ melalui
\[ [V] + [W] := [V \oplus W] \]
dan perkalian pada $R^+_k(G)$ melalui
\[ [V] \cdot [W] := [V \otimes W]. \]
Kita menulis $0$ untuk $[0]$ dan $1$ untuk $[k]$.
(Perlu diperiksa bahwa operasi-operasi ini terdefinisi dengan baik.)

Untuk representasi $U,V,W$, misalkan $u=[U]$, $v=[V]$, dan $w=[W]$.
Tuliskan $0 = [0]$ untuk representasi nol dan $1=[k]$
untuk representasi trivial.
Sekarang, isomorfisme-isomorfisme standar memberikan sifat-sifat $R^+_k(G)$ berikut:
\begin{align*}
U \oplus V &\cong V \oplus U & u+v&=v+u\\
U \oplus (V \oplus W) &\cong (U \oplus V) \oplus W & u+(v+w)&=(u+v)+w\\
0 \oplus V &\cong V & 0+v&=v\\
V \oplus 0 &\cong V & v+0&=v\\
U \otimes V &\cong V \otimes U & uv&=vu\\
U \otimes (V \otimes W) &\cong (U \otimes V) \otimes W & u(vw)&=(uv)w\\
k \otimes V &\cong V & 1v&=v\\
V \otimes k &\cong V & v1&=v\\
0 \otimes V &\cong 0 & 0v&=0\\
V \otimes 0 &\cong 0 & v0&=0\\
U \otimes (V \oplus W) &\cong (U \otimes V) \oplus (U \otimes W) &
u(v+w)&=uv+uw\\
(V \oplus W) \otimes U &\cong (V \otimes U) \oplus (V \otimes U) &
(v+w)u&=vu+wu
\end{align*}

Kita ingin menyimpulkan bahwa $R^+_k(G)$ adalah gelanggang,
tetapi invers aditif belum tersedia. Memang, karena alasan dimensi, tidak ada
representasi tak nol yang memiliki invers aditif.
Meskipun demikian, $R^+_k(G)$ memiliki struktur \emph{rig} komutatif
(``gelanggang tanpa bilangan negatif'').
Secara khusus, $R^+_k(G)$ merupakan monoid komutatif terhadap penjumlahan,
dan unsur-unsur tak nol dari $R^+_k(G)$ membentuk monoid komutatif terhadap
perkalian.

Misalkan $S_k(G)$ adalah himpunan kelas isomorfisme
representasi tak tereduksi dari $G$ atas $k$.
Menurut teorema Krull--Schmidt, setiap unsur $[V]$ dari $R^+_G(k)$
dapat ditulis secara tunggal sebagai kombinasi linear hingga
\[
[V] = m_1 [W_1] + \cdots + m_r [W_r]
\]
di mana $[W_1],\ldots,[W_r] \in S_k(G)$ dan $m_1,\ldots, m_r$ adalah
bilangan bulat tak negatif.
Dengan demikian, terdapat isomorfisme monoid aditif
$R^+_G(k) \cong \mathbb{N}^{\oplus S_k(G)}$.
(Kita akan melihat bahwa $S_k(G)$ selalu hingga, tetapi sebelum hal itu
dibuktikan, kita tetap membedakan jumlah langsung dari hasil kali langsung.)

Walaupun struktur aditif pada $R^+_G(k)$ sangat sederhana,
struktur \emph{perkaliannya} jauh lebih halus.
Misalkan $X$ adalah himpunan indeks untuk $S_k(G) = \{[W_i]\}_{i \in X}$.
Struktur perkalian ditentukan sepenuhnya oleh bilangan bulat tak negatif
$c^{ij}_\ell$ dalam persamaan
\[
[W_i] \cdot [W_j] = \sum_{\ell \in X} c^{ij}_\ell [W_\ell]
\]
ketika $i,j$ bervariasi secara bebas pada semua unsur $X$.

\begin{remark}
Bilangan $c^{ij}_\ell$ disebut \emph{bilangan Clebsch--Gordan}
untuk grup $G$. Setelah memilih basis bagi setiap $W_i$,
kita memperoleh matriks eksplisit yang mendefinisikan isomorfisme antara
$W_i \otimes W_j$ dan dekomposisinya menjadi representasi tak tereduksi;
entri-entri matriks ini disebut \emph{koefisien
Clebsch--Gordan}. Analognya untuk grup tak hingga, seperti grup ortogonal khusus
$SO(3)$, sangat penting dalam fisika kuantum dan kimia
kuantum.
\end{remark}

\subsection{Representasi Virtual}

Suatu \emph{representasi virtual} $X$ adalah unsur grup abelian bebas
yang dibangkitkan oleh $S_k(G)$.
Secara lebih konkret, terdapat suatu persamaan
\[
X = m_1 [W_1] + m_2 [W_2] + \cdots + m_r [W_r]
\]
di mana $[W_1],\ldots,[W_r]$ merupakan kelas-kelas isomorfisme yang berbeda dari
representasi tak tereduksi
dan $m_1,\ldots, m_r$ adalah bilangan bulat.
Jika semua $m_1,\ldots,m_r$ tak negatif,
maka representasi virtual tersebut adalah (kelas isomorfisme dari) suatu
representasi, menurut penafsiran yang jelas di $R^+_k(G)$.

\begin{definition}
\emph{Gelanggang representasi $R_k(G)$ dari $G$ atas $k$} adalah
grup aditif representasi virtual, dengan perkalian yang
diperoleh dengan memperluas perkalian $R^+_k(G)$ secara linear.
Jika $k=\mathbb{C}$, kita cukup menuliskannya sebagai $R(G)$.
\end{definition}

\begin{example}
Jika $G=1$ adalah grup trivial, representasi hanyalah ruang vektor.
Dalam kasus ini, $R^+_k(1) \cong \mathbb{N}$ melalui $V \mapsto \dim_k(V)$.
Gelanggang representasi virtual adalah $R_k(1) \cong \mathbb{Z}$.
\end{example}

\begin{remark}
Teori \emph{gelanggang} jauh lebih ampuh daripada teori \emph{rig}.
Oleh karena itu, kita ``menambahkan bilangan negatif'' untuk mengubah rig $R^+_k(G)$
menjadi gelanggang $R_k(G)$.
Prosedur ini merupakan analog aditif dari pembentukan gelanggang pecahan suatu gelanggang.
Secara umum, dengan cara serupa kita dapat membentuk \emph{grup Grothendieck}
atau \emph{pelengkapan grup} $G$ dari suatu monoid komutatif $M$.
Jika monoid $M$ adalah rig, grup $G$ tersebut menjadi gelanggang.
Namun, kita hanya memperoleh pembenaman $M \hookrightarrow G$ jika
$M$ bersifat \emph{kanselatif}.
\end{remark}

\section{Karakter}

Sepanjang bagian ini, kita mengasumsikan $k=\mathbb{C}$.
Untuk $z \in \mathbb{C}$, notasi $\overline{z}$ menyatakan
konjugat kompleks. Perhatikan bahwa jika $z$ adalah akar kesatuan, maka
$\overline{z}=z^{-1}$. Pengamatan sederhana ini akan membawa konsekuensi
penting.

\begin{definition}
Misalkan $G$ adalah grup hingga dan $(V,\rho)$ suatu representasi dari $G$.
Kita mendefinisikan \emph{karakter dari $\rho$} sebagai fungsi
\[
\chi_V : G \to \mathbb{C}
\]
yang diberikan oleh $\chi_V(g)=\operatorname{tr}(\rho_g)$ untuk setiap $g \in G$.
\end{definition}

Karakter memiliki banyak sifat yang berguna.


\begin{proposition}
Misalkan $(V,\rho)$ dan $(W,\sigma)$ adalah representasi dari suatu grup hingga
$G$.
Maka:
\begin{enumerate}
\item $\chi_V(1)=\dim(V)$, dengan $1$ menyatakan unsur identitas $G$,
\item $\chi_V(g^{-1}) = \overline{\chi_V(g)}$
untuk setiap $g \in G$, dan
\item $\chi_V(hgh^{-1})=\chi_V(g)$ untuk setiap $g,h \in G$.
\item
$\displaystyle
\chi_{V \oplus W}(g) = \chi_V(g) + \chi_W(g)$ untuk $g \in G$.
\item
$\displaystyle
\chi_{V \otimes W}(g) = \chi_V(g)\chi_W(g)$ untuk $g \in G$.
\item
$\displaystyle
\chi_{V^\vee}(g) =  \overline{\chi_V(g)}$ untuk $g \in G$.
\item
$\displaystyle
\chi_{\operatorname{Hom}_k(V,W)}(g) = \overline{\chi_V(g)} \chi_W(g)$ untuk $g \in G$.
\end{enumerate}
\end{proposition}


\begin{proof}
Tetapkan suatu unsur $g \in G$.
Karena $\rho(g)$ dan $\sigma(g)$ adalah matriks kompleks berorde hingga,
keduanya dapat didiagonalkan dan nilai eigennya adalah akar-akar kesatuan.
Misalkan $e_1,\ldots,e_n$ adalah basis diagonal untuk $\rho(g)$
dengan nilai eigen $\lambda_1, \ldots, \lambda_n$, dan
misalkan $f_1, \ldots, f_m$ adalah basis diagonal untuk $\sigma(g)$
dengan nilai eigen $\mu_1, \ldots, \mu_m$.
Secara khusus,
\[
\chi_V(g) = \lambda_1 + \cdots + \lambda_n \textrm{ dan }
\chi_W(g) = \mu_1 + \cdots + \mu_m.
\]
Sekarang kita buktikan setiap bagian pernyataan tersebut.

(1.)
$\chi_V(1)$ adalah matriks identitas $n \times n$, dengan $n=\dim(V)$.


(2.)
Ingat bahwa $\zeta^{-1}=\overline{\zeta}$ untuk setiap akar kesatuan
$\zeta$.
Dengan demikian,
\[
\rho_{g^{-1}}(e_i)=\lambda_i^{-1}=\overline{\lambda_i}=\overline{\rho_g(e_i)}
\]
untuk setiap $g \in G$ dan $i \in \{1,\ldots,n\}$.
Karena teras hanyalah jumlah dari semua $\lambda_i$, hasilnya
mengikuti.

(3.) Ingat bahwa $\operatorname{tr}(AB)=\operatorname{tr}(BA)$ untuk sebarang
matriks $A$ dan $B$. Jadi,
$\operatorname{tr}(ABA^{-1})=\operatorname{tr}(A^{-1}AB)=\operatorname{tr}(B)$.



(4.) Hal ini mengikuti karena $e_1, \ldots, e_n, f_1, \ldots, f_m$ merupakan
basis diagonal untuk $(\rho \oplus \sigma)(g)$ dan karena
\[
(\lambda_1 + \cdots + \lambda_n) + (\mu_1 + \cdots + \mu_m) = 
(\lambda_1 + \cdots + \lambda_n + \mu_1 + \cdots + \mu_m) \ .
\]



(5.) Perhatikan bahwa $\{ e_i \otimes f_j \}$ adalah basis diagonal untuk
$(\rho \otimes \sigma)(g)$. Hasilnya mengikuti dari pengamatan
bahwa
\[
(\lambda_1 + \cdots + \lambda_n) (\mu_1 + \cdots + \mu_m)
= \sum_{i=1}^n \sum_{j=1}^m \lambda_i \mu_j \ .
\]



(6.) Misalkan $e_1^\vee, \ldots, e_n^\vee$ menyatakan basis dual dari $e_1,
\ldots, e_n$.
Perhatikan bahwa
\[
\rho^\vee(g)(e_i^\vee)(e_j)=e_i^\vee(\rho(g^{-1})(e_j))
= \lambda_i^{-1}\delta_{ij}=\overline{\lambda_i}\delta_{ij} \ .
\]



(7.) Hal ini mengikuti dari isomorfisme
$\operatorname{Hom}_k(V,W) \cong V^\vee \otimes W$.
\end{proof}


\begin{exercise}
Gunakan notasi dari proposisi sebelumnya.
\begin{enumerate}
\item $\displaystyle
\chi_{\operatorname{Sym}^2(V)}(g)=
\chi_{\mathcal{S}^2(V)}(g)
=  \frac{1}{2}\left(\chi(g)^2 + \chi(g^2)\right)$ untuk $g \in G$.
\item $\displaystyle
\chi_{\operatorname{Alt}^2(V)}(g)=
\chi_{\Lambda^2(V)}(g)
=  \frac{1}{2}\left(\chi(g)^2 - \chi(g^2)\right)$ untuk $g \in G$.
\end{enumerate}
\end{exercise}

Proposisi di atas membuktikan bahwa semua karakter merupakan \emph{fungsi
kelas}:

\begin{definition}
Fungsi $f : G \to \mathbb{C}$ disebut \emph{fungsi kelas}
jika $f(hgh^{-1})=f(g)$ untuk semua $g,h \in G$.
Himpunan semua fungsi kelas dinotasikan dengan $C(G)$.
\end{definition}

Secara ekuivalen, $C(G)$ tepat terdiri atas fungsi-fungsi $f : G \to
\mathbb{C}$ sedemikian sehingga $f(g)=f(h)$ apabila $g, h$ berada dalam
kelas konjugasi yang sama di $G$.

Himpunan semua fungsi $f : G \to \mathbb{C}$ merupakan
aljabar-$\mathbb{C}$ terhadap penjumlahan dan perkalian titik demi titik.
Dengan kata lain, gelanggang tersebut adalah struktur gelanggang biasa pada
pangkat Kartesius $\mathbb{C}^G$.
Subhimpunan fungsi kelas $C(G)$ merupakan subaljabar gelanggang ini.
Jika $c(G)$ menyatakan banyaknya kelas konjugasi dari $G$,
maka $C(G)$ secara alami isomorfik dengan $\mathbb{C}^{c(G)}$.

\begin{example} \label{eg:regular_rep}
Misalkan $V$ adalah representasi reguler. Maka
\[
\chi_V(g) = \begin{cases}
|G| & \textrm{jika $g=1$},\\
0 & \textrm{selain itu}.
\end{cases}
\]
Ini mengikuti karena $he_g=e_g$ hanya jika $h = 1$ dalam basis alami $\{e_g\}_{g \in G}$.
\end{example}

Sifat-sifat karakter yang telah dibuktikan di atas menghasilkan
korolari berikut:

\begin{corollary}
Pemetaan suatu representasi ke karakternya
menginduksi homomorfisme gelanggang
\[
\chi_{\bullet}: R(G) \to C(G)
\]
dari gelanggang representasi ke aljabar fungsi kelas pada $G$.
\end{corollary}

Sebentar lagi kita akan melihat bahwa $\chi_{\bullet}$ bersifat \emph{injektif}.
Unsur dalam bayangan pemetaan $\chi_{\bullet}$ di atas disebut \emph{karakter
virtual}. Perhatikan bahwa $R(G)$ hanyalah gelanggang, sedangkan $C(G)$ merupakan
aljabar-$\mathbb{C}$. Keduanya tidak pernah isomorfik, tetapi kita akan melihat bahwa
$R(G) \otimes_\mathbb{Z} \mathbb{C} \cong C(G)$ sebagai aljabar-$\mathbb{C}$.

Karena himpunan kelas isomorfisme representasi tak tereduksi dari
$G$ membentuk basis-$\mathbb{Z}$ untuk $R(G)$
dan himpunan kelas konjugasi dari $G$
menginduksi basis-$\mathbb{C}$ untuk $C(G)$,
pemetaan $\chi_{\bullet}$ sepenuhnya ditentukan dengan memasangkan kedua basis ini.
Matriks yang dihasilkan disebut \emph{tabel karakter} dari $G$.

\begin{example}
Ternyata grup simetris $S_4$ memiliki tepat $5$ representasi
tak tereduksi; namai karakter-karakter tak tereduksi tersebut
sebagai $\chi_1, \ldots, \chi_5$.
Tabel karakternya adalah sebagai berikut:
\begin{center}
\begin{tabular}{|c|c|c|c|c|c|}
\hline 
 & $e $ & $(12)$ & $(12)(34)$ & $(123)$ & $(1234)$\\
\hline 
\hline 
$\chi_1$ & $1$ & $1$ & $1$ & $1$ & $1$\\
\hline 
$\chi_2$ & $1$ & $-1$ & $1$ & $1$ & $-1$\\
\hline 
$\chi_3$ & $2$ & $0$ & $2$ & $-1$ & $0$\\
\hline 
$\chi_4$ & $3$ & $1$ & $-1$ & $0$ & $-1$\\
\hline 
$\chi_5$ & $3$ & $-1$ & $-1$ & $0$ & $1$\\
\hline 
\end{tabular}
\end{center}
\end{example}

Ada beberapa konvensi umum mengenai tabel karakter.
Hampir selalu,
representasi trivial menempati baris pertama dan kelas dari
unsur identitas menempati kolom pertama.
Konvensi lain adalah mengurutkan representasi tak tereduksi menurut
dimensi yang meningkat.
Terakhir, konvensi yang lebih jarang adalah mengurutkan kelas konjugasi menurut
orde yang meningkat dari unsur mana pun yang mewakilinya.

Paket aljabar komputer dapat menghasilkan (atau mencari) tabel karakter
untuk grup hingga. Peringatan: urutan baris dan kolom pada
tabel tersebut mungkin tidak konsisten, bahkan ketika program yang persis sama dijalankan pada
data yang persis sama!

\section{Relasi Ortogonalitas}

Kita mulai dengan menunjukkan bahwa fungsi kelas menghasilkan
endomorfisme $G$-ekuivarian.

\begin{lemma} \label{lem:class_function_endo}
Misalkan $(V,\rho)$ adalah representasi kompleks dari grup hingga $G$,
misalkan $f : G \to \mathbb{C}$ adalah fungsi, dan misalkan $F \in
\operatorname{End}_k(V)$ adalah endomorfisme yang didefinisikan oleh
\[
F(v) = \sum_{g \in G} f(g)\rho_g(v) .
\]
Maka $F$ bersifat $G$-ekuivarian jika $f$ merupakan fungsi kelas.
\end{lemma}

\begin{proof}
Latihan. Ini serupa dengan argumen pengindeksan ulang dalam Teorema Maschke.
\end{proof}

Ingat bahwa jika $(V,\rho)$ adalah representasi linear dari grup hingga $G$,
maka $V^G$ adalah subruang invarian yang terdiri atas vektor
$v \in V$ sedemikian sehingga $\rho_g(v)=v$ untuk semua $g \in G$.

\begin{lemma} \label{lem:proj_onto_invariants}
Jika $(V,\rho)$ adalah representasi linear dari grup hingga $G$,
maka
\[ \dim(V^G) = \frac{1}{|G|} \sum_{g \in G} \chi_V(g). \]
\end{lemma}

\begin{proof}
Dengan argumen yang sama seperti dalam pembuktian Teorema Maschke,
transformasi linear $\pi : V \to V$ yang didefinisikan oleh
\[
\pi(v) = \frac{1}{|G|} \sum_{g \in G} \rho_g(v)
\]
merupakan proyeksi $G$-ekuivarian ke $V^G$.
Dengan memilih basis yang sesuai untuk $V$, pemetaan $\pi$ bersesuaian dengan
matriks yang memiliki $\dim(V^G)$ buah angka satu pada diagonal dan nol di tempat lain.
Perhitungan
\[
\frac{1}{|G|} \sum_{g \in G} \chi_V(g)
= \frac{1}{|G|} \sum_{g \in G} \operatorname{tr}( \rho(g) )
= \operatorname{tr}\left( \frac{1}{|G|} \sum_{g \in G}\rho(g) \right)
= \operatorname{tr}( \pi ) = \dim(V^G)
\]
menyelesaikan argumen.
\end{proof}

\begin{definition} \label{def:innerProduct}
Definisikan pasangan
\[
( \phi , \psi ) := \frac{1}{|G|}
\sum_{g \in G} \overline{\phi(g)} \psi(g)
\]
untuk sebarang dua fungsi $\phi, \psi : G \to \mathbb{C}$.
\end{definition}

\begin{proposition}
Hasil kali $(-,-)$ menjadikan $C(G)$ ruang hasil kali dalam.
\end{proposition}

\begin{proof}
Pasangan $(-,-)$ bersifat linear konjugat pada variabel pertama dan linear pada
variabel kedua. Selain itu, $(\psi,\psi)$ adalah jumlah bilangan real tak negatif
yang bernilai nol jika dan hanya jika $\psi(g)=0$ untuk setiap $g \in G$.
\end{proof}


Hasil kali dalam ini memiliki interpretasi penting:

\begin{lemma}
$\displaystyle
(\chi_V, \chi_W) = \dim_k\left( \operatorname{Hom}^G_k\left(V,
W\right)\right)$.
\end{lemma}

\begin{proof}
Jika $\tau$ adalah karakter dari $\operatorname{Hom}_k(V,W)$, maka
\[ ( \chi_V , \chi_W ) = \frac{1}{|G|}
\sum_{g \in G} \overline{\chi_V(g)} \chi_W(g)
= \frac{1}{|G|} \sum_{g \in G} \tau(g) .
\]
Dengan demikian:
\[ ( \chi_V , \chi_W ) =
\dim_k\left(\left( \operatorname{Hom}_k\left(V,W\right) \right)^G\right) =
\dim_k( \operatorname{Hom}^G_k(V,W) ) \]
seperti yang diinginkan.
\end{proof}

Suatu karakter disebut \emph{tak tereduksi} jika representasi yang bersesuaian
tak tereduksi.
Hasil utama bagian ini adalah sebagai berikut:

\begin{theorem}
Karakter-karakter tak tereduksi membentuk basis ortonormal untuk ruang
fungsi kelas $C(G)$.
\end{theorem}

\begin{proof}
Misalkan $V$ dan $W$ tak tereduksi.
Maka Lema Schur memberikan
\[
\left( \chi_V, \chi_W \right)
= \dim_k \operatorname{Hom}^G_k(V,W) =
\begin{cases}
1 & V \cong W\\
0 & V \not\cong W
\end{cases}.
\]
Jadi, karakter-karakter tak tereduksi merupakan himpunan ortonormal.
Masih perlu ditunjukkan bahwa karakter-karakter tersebut merentang $C(G)$.

Misalkan $f \in C(G)$ ortogonal terhadap setiap karakter tak tereduksi
(dan karenanya terhadap setiap karakter).
Untuk sebarang representasi $(V,\rho)$, endomorfisme
\[
F_\rho = \sum_{g \in G} \overline{f(g)} \rho_g
\]
bersifat ekuivarian menurut Lema~\ref{lem:class_function_endo}.
Dengan mengambil teras, kita memperoleh
\[
\operatorname{tr}(F_\rho) =
\sum_{g \in G} \overline{f(g)} \chi_V(g)
= |G| (f,\chi_V),
\]
yang bernilai $0$ menurut syarat ortogonalitas.
Ketika $V$ tak tereduksi,
$F_\rho$ adalah perkalian dengan skalar menurut Lema Schur;
lebih jauh, karena terasnya nol,
$F_\rho=0$.
Jika $W$ adalah subrepresentasi tak tereduksi dari $V$,
maka $F_\rho(W) \subseteq W$ karena $F_\rho$ merupakan kombinasi linear dari
fungsi-fungsi yang memiliki sifat tersebut.
Dengan demikian, kita menyimpulkan bahwa $F_\rho=0$ untuk semua representasi $\rho$.

Ketika $V$ adalah representasi reguler dengan basis $\{e_g\}_{g \in G}$,
kita memperoleh
\[
0=F_\rho(e_1)=\sum_{g \in G} \overline{f(g)} e_g .
\]
Kita harus menyimpulkan bahwa $f(g)=0$ untuk semua $g \in G$. Jadi, tidak ada fungsi kelas yang
ortogonal terhadap setiap karakter tak tereduksi.
\end{proof}

Teorema ini sangat kuat. Secara khusus, teorema ini memberitahukan bahwa
banyaknya representasi tak tereduksi bukan hanya \emph{hingga}, melainkan
juga tepat berapa jumlahnya:

\begin{corollary}
Banyaknya representasi tak tereduksi sama dengan banyaknya
kelas konjugasi dalam $G$.
\end{corollary}

Tabel karakter dapat dipandang sebagai matriks ``perubahan basis''
antara representasi tak tereduksi dalam $R(G)$ dan fungsi karakteristik
dari kelas-kelas konjugasi dalam $C(G)$.
Walaupun terdapat bijeksi antara kelas konjugasi
dan representasi tak tereduksi, kita tidak membangunnya secara eksplisit.
Memang, ada alasan kuat untuk tidak mengharapkan suatu ``bijeksi
kombinatorial'' yang eksplisit secara umum.

\begin{corollary}
Jika $V$ adalah representasi dengan dekomposisi isotipik
\[
V \cong W_1^{\oplus m_1} \oplus \cdots \oplus W_r^{\oplus m_r}
\]
untuk representasi tak tereduksi $W_1, \ldots, W_r$,
maka multiplisitasnya dapat diperoleh kembali melalui
\[
m_i = \left( \chi_V, \chi_{W_i} \right).
\]
\end{corollary}

\begin{proof}
Latihan.
\end{proof}

\begin{corollary}
Suatu karakter $\chi$ tak tereduksi jika dan hanya jika $(\chi,\chi)=1$.
\end{corollary}

\begin{proof}
Latihan.
\end{proof}

\begin{corollary}
Misalkan $G$ memiliki representasi tak tereduksi $V_1,\ldots,V_r$
dengan dimensi $n_1,\ldots, n_r$. Maka dekomposisi isotipik dari
representasi reguler adalah
\[
V_1^{n_1} \oplus \cdots \oplus V_r^{n_r} .
\]
Secara khusus, $|G|=n_1^2 + n_2^2 + \cdots + n_r^2$.
\end{corollary}

\begin{proof}
Misalkan $W$ adalah representasi reguler.
Dengan menggunakan Contoh~\ref{eg:regular_rep}, kita menghitung
\[
( \chi_{V_i} , \chi_W ) = \frac{1}{|G|}
\sum_{g \in G} \overline{\chi_{V_i}(g)} \chi_W(g)
= \frac{1}{|G|} \overline{\chi_{V_i}(1)} |G|
= \dim(V_i)=n_i.
\]
Pernyataan kedua mengikuti dengan membandingkan dimensi
representasi reguler dengan dekomposisi isotipiknya.
\end{proof}

\subsection{Menghitung Hasil Kali Dalam}

Dalam praktik, menghitung hasil kali dalam dengan menjumlahkan atas
setiap unsur grup tidaklah efisien. Kita sudah mengetahui bahwa fungsi kelas
memberikan nilai yang sama bagi setiap unsur dalam suatu kelas konjugasi.
Oleh karena itu, kita dapat menulis ulang hasil kali dalam dengan memanfaatkan
sifat ini.

\begin{proposition}
Misalkan $K_1, \ldots, K_r$ adalah kelas-kelas konjugasi dari $G$
dan misalkan $g_1,\ldots, g_r$ adalah wakil dari masing-masing kelas.
Maka
\[
( \psi, \phi ) = \sum_{i=1}^r \frac{|K_i|}{|G|}
\overline{\psi(g_i)}\phi(g_i)
\]
untuk semua fungsi kelas $\psi,\phi \in C(G)$.
\end{proposition}

Ingat bahwa transpos suatu matriks ortogonal juga ortogonal.
Rumus di atas memberikan pasangan ortonormal untuk
\emph{baris-baris} tabel karakter (setelah diboboti secara tepat menurut banyaknya unsur dalam
setiap kelas konjugasi). Kita juga memiliki ortogonalitas
\emph{kolom-kolom}, seperti berikut:

\begin{exercise}
Misalkan $\chi_1,\ldots,\chi_r$ adalah karakter-karakter tak tereduksi.
Misalkan $K(g)$ menyatakan banyaknya unsur dalam kelas konjugasi dari $g \in
G$.
Kita mempunyai
\[
\sum_{i=1}^r \overline{\chi_i(g)} \chi_i(h) =
\begin{cases}
\frac{|G|}{K(g)} & \textrm{jika $g$ berkonjugasi dengan $h$}\\
0 & \textrm{selain itu}
\end{cases}
\]
untuk $g, h \in G$.
\end{exercise}

\subsection{Dekomposisi Eksplisit}

Jika karakternya diketahui, kita dapat menghitung secara eksplisit dekomposisi
isotipik suatu representasi linear.

\begin{theorem}
Misalkan $W_1,\ldots, W_r$ adalah representasi-representasi tak tereduksi dari $G$ dengan
karakter yang bersesuaian $\chi_1,\ldots, \chi_r$.
Misalkan $(V,\rho)$ adalah suatu representasi dan $V_i$ komponen isotipik
dari $V$ yang terkait dengan $W_i$.
Maka pemetaan $\pi_i : V \to V$ yang diberikan oleh
\[
\pi_i(v) := \frac{\chi_i(1)}{|G|} \sum_{g \in G} \overline{\chi_i(g)} \rho_g(v)
\]
adalah proyeksi ke komponen isotipik $V_i$.
\end{theorem}

\begin{proof}
Dari Lema~\ref{lem:class_function_endo}, kita melihat bahwa
$\pi_i$ bersifat ekuivarian.
Jadi, $\pi_i(U) \subseteq U$ untuk setiap subrepresentasi $U$ dari $V$.
Menurut Lema Schur, $\pi_i|_U$ adalah perkalian dengan skalar apabila
$U$ tak tereduksi.
Jika $\psi$ adalah karakter suatu subrepresentasi tak tereduksi $U$ dari $V$,
maka
\[
\operatorname{tr}(\pi_i|_U)=\chi_i(1) ( \chi_i, \psi ) .
\]
Kita menyimpulkan bahwa
\[
\pi_i|_U = \frac{\chi_i(1)}{\psi(1)} ( \chi_i, \psi ) \ .
\]
Jadi, $\pi_i|_U=0$ jika $W_i \not\cong U$ dan $\pi_i|_U=\operatorname{id}$
jika $W_i \cong U$.
\end{proof}

Perhatikan bahwa teorema di atas memperumum proyeksi yang dibangun dalam
pembuktian Lema~\ref{lem:proj_onto_invariants}.

Ada pula rumus yang dapat mendekomposisi secara eksplisit suatu
representasi isotipik menjadi suku-suku langsung tak tereduksi.
Dekomposisi semacam itu jelas tidak mungkin kanonik.
Lihat \S 2.7 dari \cite{Serre}.

\section{Contoh}

\begin{example}
Misalkan $G = \mathbb{Z}/n\mathbb{Z}$. Misalkan $\zeta = e^{2\pi i/n}$
adalah akar kesatuan primitif utama berorde $n$.
Kelas-kelas konjugasi dari $G$ hanyalah unsur-unsur $G$.
Karakter-karakter tak tereduksi hanyalah
homomorfisme grup $G \to \mathbb{C}^\times$.
Karakter-karakter tersebut tepat merupakan pemetaan
\[
\chi_i(j) = \zeta^{ij}
\]
untuk $i \in 1,\ldots,n$ dan $j \in G$.
Tabel karakternya adalah sebagai berikut:
\begin{center}
\begin{tabular}{|c|c|c|c|c|c|}
\hline 
 & $0 $ & $1$ & $2$ & $\cdots$ & ${n-1}$\\
\hline 
\hline 
$\chi_0$ & $1$ & $1$ & $1$ & $\cdots$ & $1$\\
\hline 
$\chi_1$ & $1$ & $\zeta$ & $\zeta^2$ & $\cdots$ & $\zeta^{-1}$\\
\hline 
$\chi_2$ & $1$ & $\zeta^2$ & $\zeta^4$ & $\cdots$ & $\zeta^{-2}$\\
\hline 
$\vdots$ & $\vdots$ & $\vdots$ & $\vdots$ & $\ddots$ & $\vdots$\\
\hline 
$\chi_{n-1}$ & $1$ & $\zeta^{-1}$ & $\zeta^{-2}$ & $\cdots$ & $\zeta$\\
\hline
\end{tabular}
\end{center}
Ini tepat merupakan \emph{transformasi Fourier diskret}.
\end{example}

\begin{example}
Secara lebih umum, jika $G$ adalah grup abelian hingga, maka
tabel karakternya berukuran $|G| \times |G|$ karena setiap kelas konjugasi hanya memiliki
satu unsur.
Sekali lagi, karakter-karakter tak tereduksi adalah homomorfisme grup
$G \to \mathbb{C}^\times$.
Sebagai contoh, untuk grup empat Klein kita memperoleh:
\begin{center}
\begin{tabular}{|c|c|c|c|c|}
\hline 
 & $(0,0)$ & $(1,0)$ & $(0,1)$ & $(1,1)$\\
\hline 
\hline 
$\chi_{(0,0)}$ & $1$ & $1$ & $1$ & $1$\\
\hline 
$\chi_{(1,0)}$ & $1$ & $-1$ & $1$ & $-1$\\
\hline 
$\chi_{(0,1)}$ & $1$ & $1$ & $-1$ & $-1$\\
\hline 
$\chi_{(1,1)}$ & $1$ & $-1$ & $-1$ & $1$\\
\hline
\end{tabular}
\end{center}
\end{example}

\begin{example}
Misalkan $G$ adalah grup dihedral berorde $2n$:
\[
D_{2n} = \langle s,r \mid s^2, r^n, (sr)^2 \rangle
\]
dengan $n \ge 3$.

Misalkan $\zeta$ adalah akar kesatuan primitif berorde $n$.
Selain representasi trivial $\rho_0$,
selalu terdapat representasi berdimensi $1$, yaitu $\rho_1$,
dengan $\rho_1(s)=(-1)$ dan $\rho_1(r)=(1)$.
Untuk setiap bilangan bulat $1 \le i \le n-1$, kita dapat membentuk representasi berdimensi dua
$\sigma_i$ sebagai berikut:
\[
\sigma_i(s) = \begin{pmatrix} 0&1\\1&0 \end{pmatrix} \quad
\sigma_i(r) = \begin{pmatrix} \zeta^i&0\\0&\zeta^{-i} \end{pmatrix} .
\]
Namun, perhatikan bahwa $\sigma_i$ isomorfik dengan $\sigma_{n-i}$
dan, ketika $n$ genap, $\sigma_{n/2}$
tidak tak tereduksi.
Jadi, ketika $n$ genap, kita juga memiliki dua representasi
berdimensi $1$ tambahan,
$\chi_2$ dan $\chi_3$.

Ingat bahwa $\zeta^i+\zeta^{-i}=2\cos(\frac{2\pi i}{n})$.
Secara alami, tabel karakter untuk kasus genap dan ganjil perlu dibahas
secara terpisah.
Ketika $n$ ganjil, kita memperoleh:
\begin{center}
\begin{tabular}{|c|c|c|c|c|c|}
\hline 
 & $1$ & $sr^i$ & $r^{\pm 1}$ & $\cdots$ & $r^{\pm (n-1)/2}$\\
\hline 
\hline 
$\chi_0$ & $1$ & $1$ & $1$ & $\cdots$ & $1$\\
\hline 
$\chi_1$ & $1$ & $-1$ & $1$ & $\cdots$ & $1$\\
\hline 
$\sigma_{1}$ & $2$ & $0$ & $2\cos\left(\frac{2\pi}{n}\right)$ & $\cdots$
& $2\cos\left(\frac{ (n-1) \pi}{n}\right)$\\
\hline 
$\vdots$ & $\vdots$ & $\vdots$ & $\vdots$ & $\ddots$ & $\vdots$\\
\hline
$\sigma_{(n-1)/2}$ & $2$ & $0$ & $2\cos\left(\frac{(n-1)\pi}{n}\right)$
& $\cdots$ & $2\cos\left(\frac{ (n-1)^2 \pi}{n}\right)$\\
\hline
\end{tabular}
\end{center}
Ketika $n$ genap, kita memperoleh:
\begin{center}
\begin{tabular}{|c|c|c|c|c|c|c|}
\hline 
 & $1$ & $sr^{2i}$ & $sr^{2i+1}$ & $r^{\pm 1}$ & $\cdots$ & $r^{n/2}$\\
\hline 
\hline 
$\chi_0$ & $1$ & $1$ & $1$ & $1$ & $\cdots$ &  $1$\\
\hline 
$\chi_1$ & $1$ & $-1$ & $-1$ & $1$ & $\cdots$ & $1$\\
\hline 
$\chi_2$ & $1$ & $-1$ & $1$ & $-1$ & $\cdots$ &
$(-1)^{n/2}$\\
\hline 
$\chi_3$ & $1$ & $1$ & $-1$ & $-1$ & $\cdots$ & $(-1)^{n/2}$\\
\hline 
$\sigma_{1}$ & $2$ & $0$ & $0$ & $2\cos\left(\frac{2\pi}{n}\right)$ & $\cdots$
 & $-2$\\
\hline 
$\vdots$ & $\vdots$ & $\vdots$ & $\vdots$ & $\vdots$ & $\ddots$ & $\vdots$\\
\hline
$\sigma_{n/2-1}$ & $2$ & $0$ & $0$ &
$2\cos\left(\frac{(n-2)\pi}{n}\right)$ & $\cdots$ & $-2$\\
\hline
\end{tabular}
\end{center}
Kita dapat memeriksa bahwa semua representasi ini tak tereduksi dengan memastikan normanya
adalah $1$ terhadap hasil kali dalam. Dengan menjumlahkan kuadrat
nilai-nilainya pada $1$, kita melihat bahwa daftar ini lengkap.
\end{example}

\begin{example}
Tabel karakter untuk grup berganti pada $4$ huruf
adalah:
\begin{center}
\begin{tabular}{|c|c|c|c|c|}
\hline 
 & $()$ & $(1\ 2)(3\ 4)$ & $(1\ 2\ 3)$ & $(1\ 3\ 2)$\\
\hline 
\hline 
$\chi_1$ & $1$ & $1$ & $1$ & $1$\\
\hline 
$\chi_2$ & $1$ & $1$ & $\zeta$ & $\zeta^{-1}$\\
\hline 
$\chi_3$ & $1$ & $1$ & $\zeta^{-1}$ & $\zeta$\\
\hline 
$\chi_4$ & $3$ & $-1$ & $0$ & $0$\\
\hline
\end{tabular}
\end{center}
di mana $\zeta$ adalah akar kesatuan primitif berorde tiga.
\end{example}

\begin{example}
Tabel karakter untuk grup berganti pada $5$ huruf
adalah:
\begin{center}
\begin{tabular}{|c|c|c|c|c|c|}
\hline 
 & $()$ & $(1\ 2)(3\ 4)$ & $(1\ 2\ 3)$ & $(1\ 2\ 3\ 4\ 5)$ & $(1\ 3\ 5\ 2\ 4)$\\
\hline 
\hline 
$\chi_1$ & $1$ & $1$ & $1$ & $1$ & $1$\\
\hline 
$\chi_2$ & $3$ & $-1$ & $0$ & $\varphi$ & $-1/\varphi$\\
\hline 
$\chi_3$ & $3$ & $-1$ & $0$ & $-1/\varphi$ & $\varphi$\\
\hline 
$\chi_4$ & $4$ & $0$ & $1$ & $-1$ & $-1$\\
\hline 
$\chi_5$ & $5$ & $1$ & $-1$ & $0$ & $0$\\
\hline
\end{tabular}
\end{center}
di mana $\varphi$ adalah rasio emas $\frac{1}{2}(1+\sqrt{5})$.
\end{example}

%%%%%%%%%%%%%%%%%%
%%%%%%%%%%%%%%%%%%

\bibliographystyle{alpha}
\bibliography{rep_theory}

\end{document} 


